Cognitive systems must acquire information, act effectively, and preserve the resources and internal structures on which future performance depends. These demands are addressed by optimization, active inference, resource-rational analysis, and physiological regulation, but their integration remains a useful design problem. This conceptual paper proposes the Unified Variational Intelligence Framework (\uvif{}) as an organizing account of five interacting influences: epistemic pull, selective protection, computational drag, action efficiency, and sustainability. We formulate these influences as compatible vector fields subject to explicit admissibility constraints and distinguish stationary equilibria from invariant operating sets and tracking under environmental change. A restricted analytical example shows both how a stable equilibrium can arise and how the proposed formulation can coincide with constrained optimization. The intended contribution is therefore an explicit, inspectable decomposition of regulatory demands, rather than a demonstrated alternative to optimization or a universal law of intelligence. We relate the framework to cognitive theories, present hypothetical applications, and specify comparative tests involving resource manipulation, deployment horizons, irreversible outcomes, and component ablation. No empirical validation is reported; companion simulations illustrate a stipulated policy model. The framework's value depends on whether measurable, independently specified components improve prediction or system design relative to comparably resourced alternatives. Open questions concern identifiability, safety enforcement, stability under changing conditions, and the empirical justification of the five-component decomposition.
